\chapter{Contratos y garantías}
\section{Importes}
Propuesta inicial para importes y saldos: EUR, valor no negativo, dos decimales y capacidad decimal total de 12 dígitos, con máximo \code{9999999999.99}. Java empleará \code{BigDecimal}; PostgreSQL, \code{NUMERIC(12,2)}; COBOL podrá representar el mismo dominio mediante \code{PIC 9(10)V99} o doce dígitos de céntimos enteros. La elección de intercambio se cerrará antes de implementar el adaptador.

El signo contable se expresa mediante el sentido del apunte, no aceptando importes negativos arbitrarios de una transferencia. El valor monetario cero es representable; la instrucción exige un importe estrictamente positivo. Los totalizadores de lotes necesitan una capacidad propia, mayor si procede, y no reutilizarán sin análisis el límite por operación.

Se normalizan ceros decimales redundantes sin cambiar valor, como \code{1.230} a \code{1.23}. Se rechaza una entrada como \code{1.231} que necesitaría redondeo para caber en dos decimales. No se usa \code{double} como paso intermedio ni se redondea dinero de entrada silenciosamente.

\section{HTTP}
Los importes viajan como cadenas decimales con punto, por ejemplo \code{"25.00"}. Las fechas técnicas usan un formato explícito y zona horaria; los identificadores se tratan como identificadores, no como cantidades. Se distinguirán validación de petición, conflicto de idempotencia, rechazo de negocio y fallo técnico. Los endpoints, códigos HTTP y esquema de errores se definirán como OpenAPI al construir el primer caso de uso.

Una orden aceptada para procesar no equivale a una instrucción liquidada. Si el procesamiento es asíncrono, la respuesta y la consulta de estado deben expresar esa diferencia.

\section{Idempotencia}
La clave tiene un ámbito definido, propuesto como participante emisor y referencia externa. Se conserva una representación o huella canónica de los datos relevantes. Misma clave y contenido equivalente devuelve el resultado conocido o el estado en curso; misma clave con contenido diferente produce conflicto. Una restricción persistente y pruebas concurrentes respaldarán el comportamiento.

Un reintento técnico no genera una nueva identidad monetaria. Una reversión sí es una nueva operación que referencia a la original. La identidad de evento tampoco se regenera en cada intento de publicación.

\section{Java--COBOL}
Se definirán versión, codificación, terminadores, orden y tamaño de campos, signo, escala, límites, identificador de solicitud y motivos de rechazo. Los datos estructurados se separan de los logs. Los fallos técnicos deben producir código no cero; un rechazo de negocio expresado correctamente es un resultado procesado, con un estado de negocio distinto de éxito monetario.

Cada ejecución usa rutas o canales aislados. Se verifican fichero ausente, registro incompleto, tamaño excedido, caracteres inválidos, cierre y falta de respuesta. El timeout debe detener el trabajo real, no solo abandonar la espera del lanzador. No se procesará un resultado de otra ejecución por reutilizar nombres temporales.

El anterior contrato de simulación V1 de cuatro líneas no es un contrato aprobado de BatchFlow. Los formatos definitivos se diseñarán con fixtures y pruebas, conservando la distinción entre valor numérico, almacenamiento y presentación.

\section{Eventos de auditoría}
La cabecera mínima propuesta separa:
\begin{itemize}
\item Identidad y esquema: \code{eventId}, \code{eventType} y \code{schemaVersion}.
\item Entidad y orden lógico: \code{aggregateId} y \code{aggregateVersion}.
\item Contexto: \code{occurredAt}, \code{correlationId} y actor cuando corresponda.
\end{itemize}
El contenido incluye los datos mínimos para explicar la decisión, evitando secretos y datos personales innecesarios.

Los consumidores tolerarán entrega repetida y evolucionarán de forma compatible. Los nombres de topics, claves, retención, permisos y política ante eventos inválidos se cerrarán antes de introducir Kafka. La garantía de exactamente una vez de Kafka no se extiende automáticamente a una escritura en PostgreSQL.

\section{Estados y recuperación}
La instrucción debe distinguir recepción, ejecución, liquidación, rechazo y fallo técnico. El modelo definitivo aclarará qué estados son terminales y cuáles permiten reintento. Un timeout después de confirmar en base de datos no debe permitir aplicar de nuevo una operación cuyo cliente desconoce el resultado.

El lote muestra contadores y estados derivados sin ocultar instrucciones pendientes. Un fallo de publicación de auditoría no convierte una instrucción monetariamente confirmada en rechazada: se recupera la publicación desde outbox.
